\documentclass[12pt, a4paper]{article}

\usepackage[utf8]{inputenc}
\usepackage[T1]{fontenc}
\usepackage[portuguese]{babel}
\usepackage{amsmath}
\usepackage{amssymb}
\usepackage{geometry}
\usepackage{hyperref}
\usepackage{parskip}
\usepackage{titlesec}
\usepackage{booktabs}
\usepackage{dirtree}
\usepackage{helvet}
\renewcommand{\familydefault}{\sfdefault}

\geometry{margin=2.5cm}

\hypersetup{
    colorlinks=true,
    linkcolor=blue,
    urlcolor=blue
}

\title{
    \vspace{-2cm}
    \Large \textbf{IMPA Tech} \\[0.5cm]
    \large Machine Learning 2 \\
    \small Prof. Francisco Ganacim \quad Prof. Daniel Yukimura \\[1.5cm]
    \rule{\linewidth}{0.5pt} \\[0.4cm]
    {\Huge \textbf{Gaussian Football n-D}} \\[0.2cm]
    \large Documentação Técnica \\
    \rule{\linewidth}{0.5pt} \\[1.5cm]
    \normalsize
    \begin{tabular}{c}
        Gabriel Ferreira Silva \\[0.3cm]
        Letícia Aleixo Ferreira \\[0.3cm]
        Pedro Miguel Rocha Santos
    \end{tabular}
}
\date{\vspace{1cm} Rio de Janeiro, \today}

\begin{document}

\maketitle
\thispagestyle{empty}
\newpage

\tableofcontents
\newpage


% ============================================================
\section{Visão Geral do Projeto}
\label{sec:visao_geral}
% ============================================================

O sistema automatiza a geração de \textit{highlights} (a princípio, gols) de uma partida de futebol, utilizando tanto imagens quanto áudio (na forma de mel espectrogramas) extraídos de uma transmissão televisiva.

O problema é formulado sob a perspectiva de regressão. A cada clipe da partida associa-se um \textit{arousal score}, um valor escalar contínuo que representa a intensidade emocional daquele trecho. Todo clipe cujo score supera um limiar pré-definido é considerado um highlight. Para isso, a partida é subdividida em $n$ pares alinhados de clipe visual e mel espectrograma. O modelo recebe cada par como entrada e produz, como saída, um único valor escalar, o \textit{arousal score} predito para aquele clipe.

\subsection{Construção dos Labels}

Cada partida possui anotações de momentos-chave fornecidas pelo dataset SoccerNet: gols, cartões, defesas, escanteios, entre outros. Para cada gol anotado no instante $t_{\text{gol}}$, associa-se uma curva gaussiana centrada nesse instante. O \textit{arousal score} de qualquer instante $t$ da partida é definido por:

\begin{equation}
    a(t) = \exp\left(-\frac{(t - t_{\text{gol}})^2}{2\sigma^2}\right)
    \label{eq:arousal}
\end{equation}

\noindent onde:
\begin{itemize}
    \item $t_{\text{gol}}$ é o instante do gol, onde o score atinge seu pico máximo $a(t_{\text{gol}}) = 1{,}0$;
    \item $\sigma$ é calibrado de forma que o score nas bordas do intervalo
    \[
        \left[t_{\text{gol}} - \texttt{window\_before},\;
        t_{\text{gol}} + \texttt{window\_after}\right]
    \]
    seja exatamente igual ao threshold $\tau$ (default: $0{,}5$). Matematicamente,
    \[
        \sigma = \frac{d}{\sqrt{-2\ln(\tau)}}, \quad
        d = \frac{\texttt{window\_before} + \texttt{window\_after}}{2}.
    \]
    O pico da gaussiana coincide com $t_{\text{gol}}$ apenas quando
    \texttt{window\_before} $=$ \texttt{window\_after}.

    \item o score decai simetricamente à medida que $t$ se afasta de $t_{\text{gol}}$, tendendo a zero fora da janela de highlight.
\end{itemize}

\subsection{Agregação por Clipe: \texttt{np.max}}

Ao fatiar a partida em clipes, cada clipe contém múltiplos frames, cada um com um score $a(t)$ calculado pela Equação~\ref{eq:arousal}. O score final associado ao clipe é o valor máximo entre todos os seus frames:

\begin{equation}
    s_{\text{clipe}} = \max_{t \,\in\, \text{clipe}}\; a(t)
\end{equation}

A escolha de \texttt{np.max} ao invés de \texttt{np.mean} é empírica. Considere um clipe de 5 segundos a 2\,fps: esse clipe contém 10 frames. Se apenas 2 desses frames estão próximos do pico da gaussiana, a média dilui o sinal pelos 8 frames restantes, cujo score é próximo de zero. O máximo preserva o sinal mais forte do clipe, indicando que aquele clipe contém um momento de alta intensidade, independentemente de quantos frames não emocionantes coexistem nele.

\newpage

\section{Estrutura do Projeto}

O projeto está organizado em cinco diretórios principais:

\begin{itemize}
    \item \texttt{data/} — centraliza todos os dados do pipeline:
    \begin{itemize}
        \item \texttt{raw/} — arquivos brutos baixados do SoccerNet;
        \item \texttt{processed/} — clipes fatiados e mel espectrogramas, espelhando a estrutura de \texttt{raw/} no formato \texttt{liga/temporada/partida/tempo/};
        \item \texttt{labels/} — CSVs de índice e labels gerados pelos scripts de \texttt{src/}.
    \end{itemize}
    \item \texttt{models/nn/} — arquiteturas de redes neurais implementadas;
    \item \texttt{preprocessing/} — utilitários de baixo nível para fatiamento de vídeo, extração de features e geração de mel espectrogramas;
    \item \texttt{src/} — scripts principais do pipeline: download, indexação e geração de labels;
    \item \texttt{outputs/} — reservado para resultados do modelo após o treinamento.
\end{itemize}

\newpage

\subsection{Árvore de Diretórios}

\bigskip

\dirtree{%
.1 gaussian\_football/.
.2 data/.
.3 labels/.
.4 games\_index.csv.
.4 labels\_all.csv.
.3 processed/.
.3 raw/.
.2 models/.
.3 nn/.
.4 bidirectional\_lstm.py.
.4 cnnbilstm.py.
.4 vgg\_16\_lstm.py.
.2 outputs/.
.2 preprocessing/.
.3 audio\_mel\_spectogram.py.
.3 datasets\_mel\_video.py.
.3 image\_resize.py.
.3 image\_to\_array.py.
.3 loaders.py.
.3 mel\_loader.py.
.3 vgg16\_feature\_extract.py.
.3 video\_and\_audio\_get\_sequences.py.
.3 video\_color.py.
.3 video\_get\_sequence.py.
.3 video\_scorer.py.
.3 video\_slicer.py.
.2 src/.
.3 build\_index.py.
.3 build\_labels.py.
.3 download\_games.py.
.3 treinamento.ipynb.
.2 README.md.
.2 requirements.txt.
}

\newpage
\section{Pipeline de Preparação de Dados}

\subsection{Dataset e Configuração}

Os dados utilizados neste projeto são provenientes do dataset SoccerNet, que requer autorização prévia para acesso aos arquivos de vídeo. Para fins acadêmicos e não comerciais, submetemos uma solicitação de uso à equipe do SoccerNet, que nos forneceu uma senha de acesso aos arquivos protegidos.

Foram utilizados jogos da Premier League inglesa (EPL) das temporadas 2014/2015, 2015/2016 e 2016/2017, totalizando 95 partidas distribuídas da seguinte forma:


\begin{itemize}
    \item \textbf{Treino:} 58 jogos
    \item \textbf{Validação:} 19 jogos
    \item \textbf{Teste:} 18 jogos
\end{itemize}

Cada partida é composta por dois arquivos de vídeo (\texttt{1\_224p.mkv} e \texttt{2\_224p.mkv}), correspondentes ao primeiro e segundo tempo, e um arquivo de anotações (\texttt{Labels-v2.json}) são provenientes do desafio de \textit{action spotting} promovido pelo próprio SoccerNet, no qual anotadores identificaram eventos relevantes ao longo das partidas usando visão computacional.

A divisão dos jogos em treino, validação e teste segue a partição oficial do SoccerNet, definida nos arquivos \texttt{SoccerNetGames.json}, \texttt{SoccerNetGamesTrain.json}, \texttt{SoccerNetGamesValid.json} e \texttt{SoccerNetGamesTest.json}, disponíveis no repositório oficial do dataset. Esses arquivos especificam, para cada partida, o split ao qual ela pertence, garantindo que a divisão utilizada neste projeto seja reprodutível.
\subsection{Baixando os Dados — \texttt{download\_games.py}}

O script \texttt{download\_games.py} gerencia o download dos arquivos das partidas a serem utilizadas. 

A função \texttt{normalize\_splits} trata os diferentes protocolos de split do SoccerNet: a entrada pode ser um split explícito como \texttt{"train"}, um atalho como \texttt{"v1"} (equivalente a \texttt{["train", "valid", "test"]}) ou \texttt{"all"} (que inclui também o split \texttt{"challenge"}), ou ainda uma entrada inválida, que é rejeitada com uma mensagem de erro. A saída dessa função é sempre uma lista normalizada de splits válidos.

A função \texttt{game\_already\_downloaded} verifica se todos os arquivos esperados de uma partida já existem no disco. Ela é utilizada dentro de \texttt{download\_games} para garantir que jogos já baixados sejam pulados em execuções subsequentes. Por fim, o modo \texttt{--dry\_run} lista quais jogos seriam baixados e quais seriam pulados, sem efetivamente baixar nada, útil para inspecionar o estado do diretório antes de uma operação longa.

\subsection{Indexação dos Jogos — \texttt{build\_index.py}}

O script \texttt{build\_index.py} constrói o arquivo \texttt{games\_index.csv}, que funciona como um organizador central do dataset. Ele é consultado por scripts subsequentes do pipeline para saber quais jogos estão disponíveis, onde seus arquivos estão no disco e a qual split cada um pertence.

Para cada jogo, o CSV armazena as seguintes colunas:

\begin{itemize}
    \item \texttt{game\_id}: identificador legível e único do jogo, gerado pela função \texttt{make\_game\_id};
    \item \texttt{league}: código da liga (ex: \texttt{epl});
    \item \texttt{season}: temporada (ex: \texttt{2014-2015});
    \item \texttt{split}: partição do dataset (\texttt{train}, \texttt{valid} ou \texttt{test});
    \item \texttt{1\_224p.mkv}, \texttt{2\_224p.mkv}, \texttt{Labels-v2.json}: caminhos absolutos para cada arquivo da partida no disco;
    \item \texttt{valid}: booleano que indica se todos os arquivos esperados estão presentes no disco.
\end{itemize}

A função \texttt{make\_game\_id} gera o identificador a partir do caminho relativo do jogo no formato SoccerNet. A coluna \texttt{valid} é preenchida pela função \texttt{game\_already\_downloaded}, herdada do \texttt{download\_games.py}, e indica se todos os arquivos esperados estão presentes no disco para aquela partida. Jogos com \texttt{valid=False} foram identificados mas ainda não foram baixados, são excluídos do pipeline de processamento pelos scripts seguintes.

\subsection{Geração de Clipes e Labels  — \texttt{build\_labels.py}}

O script \texttt{build\_labels.py} é responsável tanto por gerar os clipes de vídeo e mel espectrograma (\texttt{generate\_clips}) quanto por construir o arquivo \texttt{labels\_all.csv} (\texttt{build\_labels}), que para cada clipe armazena:

\begin{itemize}
    \item \texttt{clip\_path}: caminho para o arquivo de vídeo do clipe;
    \item \texttt{mel\_path}: caminho para o mel espectrograma correspondente;
    \item \texttt{arousal\_score}: score de intensidade do clipe, calculado via gaussiana;
    \item \texttt{game\_id}, \texttt{season}, \texttt{split}: metadados da partida.
\end{itemize}

A função \texttt{generate\_clips} itera sobre os dois tempos de cada partida e produz, para cada intervalo gerado pelo \texttt{VideoSlicer}, um clipe de vídeo e, quando o áudio não for silencioso, um mel espectrograma. Os arquivos são organizados na estrutura \texttt{half\_1/video/} e \texttt{half\_1/mel\_spectograma/} dentro do diretório processado de cada partida. É a função \texttt{build\_labels} que realiza o cálculo do arousal score descrito na Seção ~\ref{sec:visao_geral}, consultando a gaussiana gerada pelo \texttt{VideoScorerPreprocessor} e extraindo o valor máximo de cada intervalo via \texttt{np.max}.

O \texttt{main} verifica separadamente se os clipes já foram gerados e se os labels já existem no CSV. Isso garante que em caso de interrupção, o pipeline retome do ponto correto sem regenerar dados já processados nem perder progresso parcial.

\section{Pré-processamento}

\subsection{\texttt{video\_slicer.py} — Fatiamento de Vídeo}

A classe \texttt{VideoSlicer} é responsável por fatiar vídeos em intervalos de tempo. Seu único parâmetro de inicialização é \texttt{n\_slices}, o número de clipes a serem produzidos a partir do vídeo. O método \texttt{get\_intervals} retorna uma lista de tuplas \texttt{(start, end)}, cada uma representando os instantes inicial e final de um intervalo em segundos. O usuário deve fornecer ou a duração desejada de cada intervalo (\texttt{slice\_length}), ou a quantidade desejada de intervalos (\texttt{n\_slices}), nunca os dois simultaneamente. Ao montar os intervalos, o código garante que o \texttt{end} do último clipe não ultrapasse a duração total do vídeo via \texttt{end = min(start + length, total)}.


\begin{thebibliography}{9}

\bibitem{soccernet-data}
SoccerNet.
\textit{SoccerNet Data}.
Disponível em: \url{https://www.soccer-net.org/data}.
Acesso em: \today.

\bibitem{soccernet-github}
SoccerNet.
\textit{SoccerNet GitHub Repository}.
Disponível em: \url{https://github.com/SoccerNet/SoccerNet/tree/main}.
Acesso em: \today.

\end{thebibliography}

\end{document}